% ════════════════════════════════════════════
% 模型设定与诊断（5.3.3）
% ════════════════════════════════════════════
\subsubsection{模型设定与诊断}

本节从模型设定、稳健推断、共线诊断、残差检验与预测评估五个环节建立回归分析框架，核心目标是让系数可解释，其次才比较预测精度。

模型以 log1p 曝光量为因变量，以常数项、粉丝规模对数、帖子类型、内容主题、月份、星期、时段与付费标记为解释变量，采用多元线性回归：
\begin{equation}
    y_i=\boldsymbol{x}_i^{\mathrm{T}}\boldsymbol{\beta}+\varepsilon_i,
    \qquad i=1,2,\ldots,n,
\end{equation}
其中，$\boldsymbol{\beta}$ 为待估系数，$\varepsilon_i$ 为随机扰动。模型假定条件均值为线性、扰动在给定解释变量下均值为零。系数由最小二乘估计，拟合值为解释变量与系数估计的内积，残差为观测值减拟合值。在经典假定下，系数估计的协方差等于扰动方差乘以设计矩阵交叉乘积的逆，扰动方差用残差均方估计。

曝光量重尾使残差方差随拟合值变化，经典标准误会低估不确定性。White 提出用残差平方构造协方差矩阵的稳健估计\cite{white1980}，HC3 进一步对每个观测做杠杆修正，把第 $i$ 个残差平方除以 $(1-h_{ii})^2$ 后加权，其中 $h_{ii}$ 为该观测的杠杆值。第 $j$ 个系数的标准误取协方差矩阵对角元的平方根，t 统计量为系数估计除以其标准误，置信区间按 t 分布构造。该修正不改变系数点估计，只修正标准误，因此系数表可以与经典结果直接对照。

共线性从三个角度诊断。对第 $j$ 个解释变量，以其余变量回归得到解释度 $R_j^2$，方差膨胀因子为 $1/(1-R_j^2)$；当变量以多个哑变量为一组时，改用广义方差膨胀因子，并以 $\mathrm{GVIF}^{1/(2df)}$ 与阈值 2.24 比较\cite{fox2016}。整体共线程度用标准化设计矩阵的条件数衡量，即最大奇异值与最小奇异值之比，超过 30 提示存在严重共线。为判断粉丝规模是否值得保留，另做增量 F 检验：比较不含与含粉丝规模对数的两个残差平方和，若新增变量带来的解释力相对于残差规模很小，则说明该参数没有被数据识别。

为分离共线的影响，设置五个递进模型。M0 只含常数项，作为均值基线；M1 为全模型，同时含粉丝规模对数与 11 个月份哑变量；M2 剔除粉丝规模、保留月份固定效应；M3 保留粉丝规模、剔除月份；M4 同时含粉丝规模与线性月份项。M2 用月份固定效应吸收粉丝规模的累积趋势，是本报告的最终解释模型；M1、M3、M4 用于检验付费与内容类型的系数是否随设定改变。

对最终模型依次做四类残差检验。正态性用 Jarque--Bera 统计量，它由残差偏度与超额峰度构造，服从自由度为 2 的卡方分布\cite{jarque1980}。异方差用 Breusch--Pagan 辅助回归，将残差平方对解释变量回归，统计量为样本量乘以辅助回归的解释度\cite{breusch1979}；White 检验则不预设异方差的具体形式，直接用残差平方对解释变量的平方与交叉项回归。函数形式用 Ramsey RESET 检验，在回归中加入拟合值的平方与立方后做联合 F 检验\cite{ramsey1969}。强影响点用 Cook 距离衡量，它同时考虑残差大小与杠杆值，超过 1 视为强影响点；另报告 Durbin--Watson 统计量，作为残差序列相关的参考。

为判断样本内拟合是否过拟合，采用十折交叉验证，重复 5 次并固定随机种子 20261122。每次把样本随机分为十份，用九份训练、一份测试，记录测试折上的拟合优度、均方根误差与平均绝对误差；拟合优度以测试折自身均值为基准，可以为负，负值表示预测还不如直接用测试折均值。各项指标在 50 个测试折上取平均。

为回答“曝光方差中有多少可归因于发帖决策”，把 7 个变量组记为集合，对全部 $2^7=128$ 个子集分别计算回归的解释度，再对每个变量组在所有加入顺序上的边际贡献取平均，得到 Shapley 贡献。各组贡献之和等于全模型解释度，未解释部分为 1 减去该解释度。与逐次加入变量的顺序分解相比，Shapley 分解不受变量排列影响，适合处理月份与粉丝规模这类高度相关的组\cite{shapley1953}。

综上所述，本节建立了由回归估计、稳健推断、共线诊断、残差检验、交叉验证与方差分解组成的完整框架，得到了可复核的模型选择规则，接下来将基于此给出求解结果。
